
\begingroup
\shorthandoff{<>}% Babel español trata de otro modo los delimitadores de curvatura de Xy-pic.
\begin{proposition}
\label{XIII.1.13}
Sea $f \colon X \to Y$ un $S$-morfismo.
\begin{enumerate}
\item[1)] Sea
$$
\xymatrix@C=.5cm{ F \ar[r] & G \ar@<2pt>[r] \ar@<-2pt>[r] & H }
$$
un diagrama exacto de haces de conjuntos sobre~$X$. Supongamos que
$(G,f)$ es cohomológicamente propio relativamente a~$S$ en dimensión
$\leqslant 0$ y que $(H,f)$ es cohomológicamente propio relativamente
a~$S$ en dimensión $\leqslant -1$. Entonces $(F,f)$ es
cohomológicamente propio relativamente a~$S$ en dimensión
$\leqslant 0$.

\item[2)] Sea
\marginpar{363}
$F \to G$ un monomorfismo de haces de grupos sobre~$X$. Si $Y_1$ es un
esquema étale sobre~$Y$, se pone $X_1 = Y_1 \times_Y X$, y se denota por
$f_1$ (\resp $F_1$, \resp $G_1$) la imagen inversa de~$f$ (\resp $F$,
\resp $G$) sobre~$Y_1$ (\cf \Ref{eq:XIII.1.0.2}). Supongamos que
$(G,f)$ es cohomológicamente propio relativamente a~$S$ en dimensión
$\leqslant 0$ (\resp en dimensión $\leqslant 1$) y que, para todo
esquema $Y_1$ étale sobre~$Y$ y todo torsor $Q$ bajo $G_1$,
$(Q/F_1,f_1)$ es cohomológicamente propio relativamente a~$S$ en
dimensión $\leqslant -1$ (\resp en dimensión $\leqslant 0$). Entonces
$(F,f)$ es cohomológicamente propio relativamente a~$S$ en dimensión
$\leqslant 0$ (\resp en dimensión $\leqslant 1$).

\item[3)] Sea $F \to G$ un monomorfismo de haces de grupos sobre~$X$.
Supongamos que $(F,f)$ es cohomológicamente propio relativamente a~$S$
en dimensión $\leqslant 1$ y que $(G,f)$ es cohomológicamente propio
relativamente a~$S$ en dimensión $\leqslant 0$. Entonces, para todo
torsor $Q$ bajo~$G$, $(Q/F,f)$ es cohomológicamente propio relativamente
a~$S$ en dimensión $\leqslant 0$.
\end{enumerate}
\end{proposition}

\subsubsection*{Demostración}
1) Sea $\Phi$ (\resp $\Phi_1$, \resp $\Phi_2$) el stack en categorías
discretas asociado a~$F$ (\resp $G$, \resp $H$), y sea $\Phi_3$ el stack
final sobre~$X$. Se tiene entonces un diagrama exacto
$$
\xymatrix@C=.5cm{ \Phi \ar[r] & \Phi_1 \ar@<2pt>[r] \ar@<-2pt>[r] &
\Phi_2 \ar@<4pt>[r] \ar[r] \ar@<-4pt>[r] & \Phi_3 }
\text{.}
$$
Por hipótesis, $(\Phi_1,f)$ es cohomológicamente propio relativamente
a~$S$ en dimensión $\leqslant 1$ y $(\Phi_2,f)$ es cohomológicamente
propio relativamente a~$S$ en dimensión $\leqslant 0$
\eqref{XIII.1.2}. Como $(\Phi_3,f)$ es evidentemente cohomológicamente
propio relativamente a~$S$ en dimensión $\leqslant -1$, se deduce
de~\Ref{XIII.1.12} que $(\Phi,f)$ es cohomológicamente propio
relativamente a~$S$ en dimensión $\leqslant 1$, \ie que $(F,f)$ es
cohomológicamente propio relativamente a~$S$ en dimensión
$\leqslant 0$.

\smallskip
2) Demostremos primero que, si $(G,f)$ es cohomológicamente propio
relativamente a~$S$ en dimensión $\leqslant 0$ y si los
$(Q/F_1,f_1)$ son cohomológicamente propios relativamente a~$S$ en
\marginpar{364}
dimensión $\leqslant -1$, entonces $(F,f)$ es cohomológicamente propio
relativamente a~$S$ en dimensión $\leqslant 0$. Según
\Ref{XIII.1.3.1}, basta demostrar que, para todo esquema $Y_1$ étale
sobre~$Y$ y todo torsor $P$ sobre~$X_1$ bajo~$F_1$,
$(\leftexp{P\mkern-4mu}{F}_1,f_1)$ es cohomológicamente propio
relativamente a~$S$ en dimensión $\leqslant 0$ cuando se considera
$\leftexp{P\mkern-4mu}{F}_1$ como haz de conjuntos, y que el morfismo
canónico
$$
d \colon g^*(\R^1 f_* F) \to \R^1 f'_* F'
$$
es inyectivo. La primera afirmación se deduce inmediatamente de 1),
pues, si $Q$ denota el torsor obtenido de~$P$ por extensión del grupo
estructural $F_1 \to G_1$, se tiene un isomorfismo
$\leftexp{Q}{G}_1 / \leftexp{P\mkern-4mu}{F}_1 \isomto Q/F_1$.

Demostremos que $d$ es inyectivo. Basta demostrar que, si $Y_1$ es un
esquema étale sobre~$Y$, si $P$ y $\tilde{P}$ son dos torsores bajo
$F_1$ cuyas imágenes inversas $P'$ y~$\tilde{P}'$ sobre~$X'_1$ son
isomorfas, entonces, tras efectuar una extensión étale sobreyectiva
de~$Y_1$, $P$ y~$\tilde{P}$ se vuelven isomorfos. Elijamos un
isomorfismo $p' \colon P' \isomto \tilde{P}'$. Sean $Q$
(\resp~$\tilde{Q}$) el torsor obtenido de~$P$ (\resp $\tilde{P}$) por
extensión del grupo estructural $F_1 \to G_1$. Las imágenes inversas
$Q'$ (\resp $\tilde{Q}'$) de~$Q$ (\resp $\tilde{Q}$) sobre~$X'_1$ se
obtienen de~$P'$ (\resp $\tilde{P}'$) por
\ifthenelse{\boolean{orig}}
{extensión}
{la extensión}
del grupo estructural $F'_1 \to G'_1$; sea
$q' \colon Q' \isomto \tilde{Q}'$ el isomorfismo que se obtiene del
mismo modo a partir de~$p'$. Como $(G,f)$ es cohomológicamente propio
relativamente a~$S$ en dimensión $\leqslant 0$, puede suponerse, tras
efectuar una extensión étale sobreyectiva de~$Y_1$, que $q'$ es la
imagen de un isomorfismo $q \colon Q \isomto \tilde{Q}$. Al torsor~$P$
(\resp~$\tilde{P}$) se asocia una sección $x$ de~$Q/F_1$ (\resp una
sección $\tilde{x}$ de~$\tilde{Q}/F_1$), y, para que $P$ y
$\tilde{P}$ sean isomorfos, es necesario y suficiente que exista un
isomorfismo $Q \isomto \tilde{Q}$ tal que el isomorfismo
$$
e \colon \H^0(X_1,Q/F_1) \to \H^0(X_1,\tilde{Q}/F_1)
$$
que induce lleve $x$ a~$\tilde{x}$. Tomamos el isomorfismo~$q$. Las
secciones $e(x)$ y $\tilde{x}$ de~$\H^0(X_1,\tilde{Q}/F_1)$ tienen la
misma imagen en~$\H^0(X'_1,\tilde{Q}'/F'_1)$. Como
$(\tilde{Q}/F_1,f_1)$ es cohomológicamente propio relativamente a~$S$
en dimensión $\leqslant -1$,
\marginpar{365}
tras efectuar una extensión étale sobreyectiva de~$Y_1$, se tiene
efectivamente $e(x)=\tilde{x}$, lo que demuestra la inyectividad de~$d$.

Para concluir la demostración, queda demostrar que, si $(G,f)$ es
cohomológicamente propio relativamente a~$S$ en dimensión
$\leqslant 1$, y si, para todo esquema $Y_1$ étale sobre~$Y$ y todo
torsor $Q$ sobre~$X_1$ bajo
\ifthenelse{\boolean{orig}}
{$F_1$}
{$G_1$},
$(Q/F_1,f_1)$ es cohomológicamente propio relativamente a~$S$ en
dimensión $\leqslant 0$, entonces el morfismo $d$ es sobreyectivo. Sea
$P'$ un torsor sobre~$X'_1$ bajo~$F'_1$, y sea $Q'$ el torsor bajo
$G'_1$ obtenido a partir de~$P'$ por extensión del grupo estructural.
Dar~$P'$ equivale a dar~$Q'$ y una sección
$x'$ de~$\H^0(X'_1,Q'/F'_1)$. Se deduce entonces de la sobreyectividad
del morfismo
$$
g^*(\R^1 f_* G) \to \R^1 f'_* G'
$$
que, tras efectuar una extensión étale sobreyectiva de~$Y_1$, existe
un torsor $Q$ bajo~$G_1$ cuya imagen inversa sobre~$X'_1$ es isomorfa
a~$Q'$. Utilizando que $(Q/F_1,f_1)$ es cohomológicamente propio
relativamente a~$S$ en dimensión $\leqslant 0$, puede suponerse
asimismo que existe un elemento $x$ de~$\H^0(X_1,Q/F_1)$ cuya imagen
en~$\H^0(X'_1,Q'/F'_1)$ es~$x'$. Los datos de~$Q$ y~$x$ determinan un
torsor $P$ bajo~$F_1$, cuya imagen inversa sobre
\ifthenelse{\boolean{orig}}
{$X_1$}
{$X'_1$}
es isomorfa a~$P'$, lo que demuestra la sobreyectividad de~$d$.

\smallskip
3) Demostremos que $(Q/F,f)$ es cohomológicamente propio relativamente
a~$S$ en dimensión $\leqslant -1$, \ie que, para todo $S$-esquema
$S'$, para todo esquema $Y_1$ étale sobre~$Y$, si $x$ y~$\tilde{x}$
son dos elementos de~$\H^0(X_1,Q_1/F_1)$ cuyas imágenes $x'$ y
$\tilde{x}'$ en~$\H^0(X'_1,Q'_1/F'_1)$ son iguales, entonces, tras
una extensión sobreyectiva de~$Y_1$, se tiene $x=\tilde{x}$. A~$x$
(\resp~$\tilde{x}$) se asocia un torsor $P$ (\resp~$\tilde{P}$) bajo
$F_1$, tal que $Q_1$ se obtiene de~$P$ (\resp~$\tilde{P}$) por la
extensión $F_1 \to G_1$ del grupo estructural. De la igualdad
$x'=\tilde{x}'$ se deduce que existe un isomorfismo
$u' \colon P' \isomto \tilde{P}'$ tal que el isomorfismo inducido
sobre~$Q'_1$ por~$u'$ es la identidad. Como $(F,f)$ es
cohomológicamente propio relativamente a~$S$ en dimensión
$\leqslant 0$, se deduce que,
\marginpar{366}
tras una extensión étale sobreyectiva de~$Y_1$, existe un isomorfismo
$u \colon P \to \tilde{P}$ que levanta~$u'$; el hecho de que $(G,f)$
sea cohomológicamente propio relativamente a~$S$ en dimensión
$\leqslant -1$ implica entonces que $x=\tilde{x}$.

Demostremos que $(Q/F,f)$ es cohomológicamente propio relativamente
a~$S$ en dimensión $\leqslant 0$. Sean $Y''$ un esquema étale sobre
$Y'$ y $x''$ un elemento de~$\H^0(X'',Q''/F'')$. A~$x''$ se asocia un
torsor $P''$ sobre~$X''$ bajo~$F''$. Como $(F,f)$ es
cohomológicamente propio relativamente a~$S$ en dimensión
$\leqslant 1$, pueden encontrarse morfismos étale sobreyectivos
$Y''_1 \to Y''$ y $Y_1 \to Y$, tales que existe un morfismo
$Y''_1 \to Y'_1$, y un torsor $P$ sobre~$X_1$ bajo~$F_1$ cuya imagen
inversa sobre~$X''_1$ es isomorfa a la imagen inversa de~$P''$. Del
hecho de que $(G,f)$ es cohomológicamente propio relativamente a~$S$
en dimensión $\leqslant 0$ se deduce entonces que pueden elegirse
$Y''_1$ y $Y_1$ de modo que el torsor obtenido de~$P$ por extensión
del grupo estructural $F_1 \to G_1$ sea isomorfo a~$Q_1$; a~$P$
corresponde un elemento $x$ de~$\H^0(X_1,Q_1/F_1)$, cuya imagen en
$\H^0(X''_1,Q''_1/F''_1)$
\ifthenelse{\boolean{orig}}
{es isomorfa a}
{es igual a}
la imagen inversa de~$x''$,
lo que concluye la demostración.
\endgroup
